%==============================================================================
\section{Support Vector Machine (SVM)}
\label{sec:svm}
%==============================================================================

\begin{dinhnghia}[title={Support Vector Machine (SVM)}]
Thuật toán học có giám sát mạnh và linh hoạt, dùng cho phân loại và hồi quy; thường dùng cho \textbf{phân loại}.
Ra đời từ thập niên 1960, hoàn thiện thập niên 1990; phổ biến nhờ xử lý nhiều biến liên tục và phân loại.
\end{dinhnghia}

\subsection{Cách SVM hoạt động}

\begin{dongco}[title={Mục tiêu}]
Tìm \textbf{siêu phẳng} (hyperplane) tách các lớp; tối đa hóa \textbf{margin} --- khoảng cách từ siêu phẳng đến các điểm gần nhất mỗi lớp (support vectors).
\end{dongco}

\begin{phuongphap}[title={Khoảng cách đến siêu phẳng}]
\[
\mathrm{distance} = \frac{w \cdot x + b}{\|w\|}
\]
$w$: vector trọng số (vuông góc siêu phẳng); $b$: bias; $\|w\|$: chuẩn Euclid của $w$.
\end{phuongphap}

\begin{phuongphap}[title={Bài toán tối ưu}]
Tối đa hóa margin với ràng buộc mọi điểm được phân loại đúng --- bài toán lồi, giải bằng quadratic programming.
Nếu không tách tuyến tính được: dùng \textbf{kernel trick} ánh xạ lên không gian chiều cao hơn mà không tính ánh xạ tường minh.
\end{phuongphap}

\tsimg{06_svm/img01.jpg}

\begin{dinhnghia}[title={Support Vectors}]
Điểm dữ liệu \textbf{gần siêu phẳng nhất}; định nghĩa đường phân tách.
\end{dinhnghia}

\begin{dinhnghia}[title={Hyperplane}]
Mặt phẳng quyết định chia không gian thành vùng lớp khác nhau (đường thẳng trong 2D).
\end{dinhnghia}

\begin{dinhnghia}[title={Margin}]
Khoảng cách giữa hai đường song song qua support vectors gần nhất hai lớp; margin \textbf{lớn} thường tốt hơn margin nhỏ.
\end{dinhnghia}

\subsection{Triển khai SVM bằng Python}

\begin{vidu}[title={Import và dữ liệu mẫu}]
\begin{lstlisting}
import numpy as np
import matplotlib.pyplot as plt
from sklearn.datasets import make_blobs

X, y = make_blobs(n_samples=100, centers=2, random_state=0,
                  cluster_std=0.50)
plt.scatter(X[:, 0], X[:, 1], c=y, s=50, cmap='summer')
\end{lstlisting}
\tsimg{06_svm/img02.jpg}
\end{vidu}

\begin{vidu}[title={Nhiều đường phân tách có thể}]
\begin{lstlisting}
xfit = np.linspace(-1, 3.5)
plt.scatter(X[:, 0], X[:, 1], c=y, s=50, cmap='summer')
for m, b in [(1, 0.65), (0.5, 1.6), (-0.2, 2.9)]:
    plt.plot(xfit, m * xfit + b, '-k')
plt.xlim(-1, 3.5)
\end{lstlisting}
\tsimg{06_svm/img03.jpg}
\end{vidu}

\begin{vidu}[title={Maximum Marginal Hyperplane (MMH)}]
\begin{lstlisting}
for m, b, d in [(1, 0.65, 0.33), (0.5, 1.6, 0.55), (-0.2, 2.9, 0.2)]:
    yfit = m * xfit + b
    plt.plot(xfit, yfit, '-k')
    plt.fill_between(xfit, yfit - d, yfit + d, color='#AAAAAA', alpha=0.4)
plt.xlim(-1, 3.5)
\end{lstlisting}
\tsimg{06_svm/img04.jpg}

SVM chọn đường \textbf{tối đa hóa margin}.
\end{vidu}

\begin{vidu}[title={SVC kernel tuyến tính}]
\begin{lstlisting}
from sklearn.svm import SVC
model = SVC(kernel='linear', C=1E10)
model.fit(X, y)
\end{lstlisting}
\end{vidu}

\begin{vidu}[title={Hàm decision\_function và support vectors}]
Vẽ biên quyết định và support vectors (xem code đầy đủ trong notebook); sau khi fit:

\begin{lstlisting}
print(model.support_vectors_)
\end{lstlisting}

\textbf{Output mẫu:}
\begin{verbatim}
[[0.5323772  3.31338909]
 [2.11114739 3.57660449]
 [1.46870582 1.86947425]]
\end{verbatim}
\tsimg{06_svm/img05.jpg}
\end{vidu}

\subsection{SVM Kernels}

\begin{phuongphap}[title={Kernel trick}]
Kernel biến không gian đầu vào chiều thấp sang dạng phù hợp; biến bài toán không tách tuyến tính thành tách được trong không gian mới.
\end{phuongphap}

\begin{tinhchat}[title={Linear kernel}]
$k(x, x_i) = \sum x \cdot x_i$ (tích vô hướng).
\end{tinhchat}

\begin{tinhchat}[title={Polynomial kernel}]
$K(x, x_i) = 1 + \sum(x \cdot x_i)^d$; $d$ là bậc đa thức (chỉ định thủ công).
\end{tinhchat}

\begin{tinhchat}[title={RBF kernel}]
$K(x, x_i) = \exp(-\gamma \sum (x - x_i)^2)$; $\gamma \in [0,1]$, mặc định thường 0.1.
\end{tinhchat}

\begin{vidu}[title={SVM với kernel trên Iris (2 feature)}]
\begin{lstlisting}
from sklearn import svm, datasets
import matplotlib.pyplot as plt
import numpy as np

iris = datasets.load_iris()
X = iris.data[:, :2]
y = iris.target

x_min, x_max = X[:, 0].min() - 1, X[:, 0].max() + 1
y_min, y_max = X[:, 1].min() - 1, X[:, 1].max() + 1
h = (x_max - x_min) / 100
xx, yy = np.meshgrid(np.arange(x_min, x_max, h),
                     np.arange(y_min, y_max, h))
X_plot = np.c_[xx.ravel(), yy.ravel()]
C = 1.0

svc = svm.SVC(kernel='linear', C=C).fit(X, y)
Z = svc.predict(X_plot).reshape(xx.shape)
plt.contourf(xx, yy, Z, alpha=0.3)
plt.scatter(X[:, 0], X[:, 1], c=y)
plt.title('SVC with linear kernel')
\end{lstlisting}
\tsimg{06_svm/img06.jpg}

Kernel RBF:
\begin{lstlisting}
svc = svm.SVC(kernel='rbf', gamma='auto', C=C).fit(X, y)
\end{lstlisting}
\tsimg{06_svm/img07.jpg}
\end{vidu}

\subsection{Tinh chỉnh tham số SVM}

\begin{phuongphap}[title={Tham số quan trọng}]
\textbf{kernel}: linear, polynomial, RBF, sigmoid.
\textbf{C}: cân bằng margin vs lỗi phân loại; $C$ lớn $\Rightarrow$ ít sai số hơn, margin có thể nhỏ hơn.
Tham số riêng từng kernel: bậc đa thức, $\gamma$ RBF, v.v.
Dùng \textbf{cross-validation} (GridSearchCV) để chọn bộ tham số tốt.
\end{phuongphap}

\begin{vidu}[title={GridSearchCV}]
\begin{lstlisting}
from sklearn.model_selection import GridSearchCV
from sklearn.svm import SVC

param_grid = {
    'C': [0.1, 1, 10, 100],
    'kernel': ['linear', 'poly', 'rbf', 'sigmoid'],
    'degree': [2, 3, 4],
    'gamma': ['scale', 'auto']
}
grid_search = GridSearchCV(SVC(), param_grid, cv=5)
grid_search.fit(X_train, y_train)
print("Best parameters:", grid_search.best_params_)
print("Best accuracy:", grid_search.best_score_)
\end{lstlisting}

\textbf{Output mẫu:}
\begin{verbatim}
Best parameters: {'C': 0.1, 'degree': 3, 'gamma': 'scale', 'kernel': 'poly'}
Best accuracy: 0.975
\end{verbatim}
\end{vidu}

\subsection{Ưu và nhược điểm}

\begin{tinhchat}[title={Ưu điểm}]
Độ chính xác cao; hiệu quả không gian chiều cao; dùng subset điểm huấn luyện nên tiết kiệm bộ nhớ.
\end{tinhchat}

\begin{luuy}[title={Nhược điểm}]
Thời gian huấn luyện lớn --- không phù hợp dataset rất lớn; kém khi các lớp chồng lấn nhiều.
\end{luuy}
